\documentclass[12pt]{article}
\usepackage[T1]{fontenc}
\usepackage[utf8]{inputenc}
\usepackage{lmodern}
\usepackage{textcomp}
\usepackage{enumitem}
\usepackage{geometry}
\geometry{margin=1in}
\begin{document}
\begin{center}
Answer Key - Constitutional Law - Graduate Level Assessment 9 \
\small (Answers are ordered exactly as in the question paper.)
\end{center}

\section*{Multiple Choice}
\begin{enumerate}[label=\arabic*.]
\item \textbf{Answer:} C
\\ \textit{Options:} A) The proceedings are conducted in secret.; B) There is no right to counsel.; C) There is a right to Miranda warnings.; D) There is no right to have evidence excluded.
\\ \textit{Note:} The correct answer is accurate because individuals testifying before a grand jury do not have the right to Miranda warnings, as these warnings are only required during custodial interrogations by law enforcement, not in the context of a grand jury proceeding where witnesses are not considered in custody.
\item \textbf{Answer:} D
\\ \textit{Options:} A) The privileges and immunities clause of the Fourteenth Amendment.; B) The comity clause of Article IV.; C) The Fifth Amendment's prohibition against compulsory self-incrimination.; D) The equal protection clause of the Fourteenth Amendment.
\\ \textit{Note:} The equal protection clause of the Fourteenth Amendment applies to corporations because the Supreme Court has recognized that corporations are considered "persons" under the law, thereby granting them certain constitutional rights, including the right to equal protection under the law.
\item \textbf{Answer:} A
\\ \textit{Options:} A) Must be separately signed if the offeree supplies a form contract containing the offer.; B) Is valid for three months.; C) Is nonassignable.; D) Can not exceed a three-month duration even if consideration is given.
\\ \textit{Note:} The correct answer is based on the principle that an irrevocable written offer, or firm offer, requires the offeree's acceptance to be unequivocal and must be separately signed if the offeree presents a form contract that includes the offer, ensuring clarity and mutual agreement in the contract formation process.
\item \textbf{Answer:} A
\\ \textit{Options:} A) It is an agreement between two or more to commit a crime.; B) It is solicitation of another to commit a felony.; C) No act is needed other than the solicitation.; D) Withdrawal is generally not a defense.
\\ \textit{Note:} The correct answer is that solicitation does not involve an agreement between two or more parties to commit a crime, as solicitation is defined as the act of enticing or urging another person to commit a crime rather than a collaborative agreement to engage in criminal activity.
\item \textbf{Answer:} C
\\ \textit{Options:} A) A private individual from discriminating against a person based on race.; B) A private individual from discriminating against a person based on nationality.; C) A state official from discriminating against a person based on race.; D) A federal official from discriminating against a person based on nationality.
\\ \textit{Note:} The correct answer is supported by the enabling clause of the Fourteenth Amendment, which grants states the authority to enforce laws against discrimination, thus allowing them to regulate state officials' actions to prevent racial discrimination.
\end{enumerate}

\section*{True / False}
\begin{enumerate}[label=\arabic*.]
\setcounter{enumi}{5}
\item \textbf{Answer:} True
\\ \textit{Options:} A) True; B) False
\\ \textit{Note:} In most states, the classification of homicide into degrees-such as first-degree, second-degree, or manslaughter-is largely based on the defendant's mental state, including intent, premeditation, and recklessness, at the time of the offense.
\item \textbf{Answer:} True
\\ \textit{Options:} A) True; B) False
\\ \textit{Note:} A contract can be discharged for frustration of purpose when an unforeseen event fundamentally changes the nature of the contract's underlying purpose, regardless of whether the event was foreseeable at the time of formation.
\item \textbf{Answer:} False
\\ \textit{Options:} A) True; B) False
\\ \textit{Note:} The statement is false because a warranty of merchantability ensures that a product is fit for ordinary use and meets general quality standards, but it does not guarantee that the product will meet the buyer's specific needs, which may require seller expertise or input.
\item \textbf{Answer:} True
\\ \textit{Options:} A) True; B) False
\\ \textit{Note:} Reversions are a legal mechanism in property law whereby the original owner of a property regains ownership rights automatically when a temporary estate, such as a life estate or term of years, ends, thus confirming the statement as true.
\item \textbf{Answer:} True
\\ \textit{Options:} A) True; B) False
\\ \textit{Note:} An implied-in-fact contract is established when a client accepts an attorney's services without a fee agreement because the actions of both parties indicate a mutual understanding that the attorney's services are being rendered with the expectation of compensation, even in the absence of a formal written contract.
\end{enumerate}

\section*{Long Form Response}
\begin{enumerate}[label=\arabic*.]
\setcounter{enumi}{10}
\item \textbf{Answer:} The President's power to appoint officials with the advice and consent of the Senate plays a crucial role in maintaining the balance of power within the federal government. This system of checks and balances ensures that while the President can nominate individuals to key positions, such as federal judges and cabinet members, the Senate retains the authority to scrutinize and confirm these appointments. This collaborative process prevents any single branch from wielding excessive power and fosters accountability and transparency in governance. Additionally, it encourages a dialogue between the executive and legislative branches, reflecting a broader democratic principle that no one entity should dominate the federal administration. Ultimately, this dynamic reinforces the constitutional framework designed to limit government overreach and protect civil liberties.
\\ \textit{Note:} correct because it emphasizes the collaborative nature of the appointment process between the President and the Senate, illustrating how it serves as a check on executive power and reinforces the principle of checks and balances within the federal government.
\item \textbf{Answer:} The warranty of merchantability is a fundamental principle in consumer protection law that ensures goods sold by merchants meet a minimum standard of quality and are fit for their intended purpose. When a buyer purchases a defective item, this warranty entitles them to seek remedies, such as repair, replacement, or refund, thereby safeguarding their rights against substandard products. Under constitutional law, this warranty supports consumer rights by promoting fair trade practices and preventing deceptive practices that could undermine the economic well-being of individuals. As such, it not only establishes a legal expectation for merchants but also reinforces the broader constitutional principles of fairness and justice in commercial transactions, ensuring that consumers can rely on the quality of the goods they purchase.
\\ \textit{Note:} The answer correctly identifies the warranty of merchantability as a crucial consumer protection mechanism that ensures goods meet basic quality standards, thereby enabling buyers to seek remedies for defective items and reinforcing consumer rights under constitutional law.
\end{enumerate}

\vfill
\noindent\textit{End of Answer Key}
\end{document}